\section{Related Work}
\subsection{Model Repair in Symbolic Planning}
Creative problem solving can be described as a failure in an agent's
understanding of the world such that it is insufficient to complete a
given task \cite{ijcai2023p777}. Existing literature has explored
various approaches to this problem. In one line of work reinforcement
learning (RL) is used to discover action primitives after reaching an
unrecoverable state \cite{rapid_learn, lorang_rsj_2024}. \citet{gizzi_ieee_2019}
propose a system that enumerates action-object combinations after
reaching a failure case. Rather than learning
entirely new action primitives or having to test many different
action-object pairs, our approach seeks to use an LLM to offer object
substitions.

Repair approaches address planning tasks that are unsolvable as specified, asking how the specification can be modified so that a solution exists~\cite{bercher2025modelrepair}. \citet{gragera2023repair}, for example, perform repair by fixing incomplete action effects by compiling the unsolvable task into an extended task that may insert the missing effects. Our approach instead works by reassigning the types of objects already present in the problem, leaving the domain, goal, and initial state intact. Most relevant to our design, \citet{gragera2023exploring} prompt an LLM \emph{stand-alone} to fix unsolvable tasks and find its limited reasoning insufficient for identifying repairs that yield solvable tasks. We reach a different outcome by changing the LLM's role: rather than trusting it to repair the task, we confine it to proposing type reassignments and delegate all verification to the symbolic planner.

\subsection{LLMs for Planning}
The use of LLMs for symbolic planning has been widely explored, with
mixed results. \citet{motwani2026longcot} find that LLMs struggle with long-horizon Chain of Thought (CoT) tasks in general. Benchmarking that uses LLMs for plan generation confirms that LLMs do not reliably produce correct plans
\citep{valmeekam_large_nodate, valmeekam2023planbenchextensiblebenchmarkevaluating}. Despite these limitations, LLMs have shown promise as components within larger planning pipelines rather than as end-to-end planners. \citet{liu_llmp_2023} use LLMs to generate PDDL problem files from natural language descriptions, while \citet{singh_twostep_2024} demonstrate that LLMs can decompose problems into simpler sub-tasks amenable to parallel execution. Researchers have similarly integrated LLMs into Hierarchical Task Networks to break high-level goals into more tractable sub-tasks \citep{pmlr-v288-munoz-avila25a, xu2025online}. LLMs have also been used to find symbolic models given pre-existing black-box style executors \cite{yang2026skillwrappergenerativepredicateinvention}.

Another line of research uses LLM-generated outputs to guide or improve
search and learning. \citet{silver2022pddl} propose treating an LLM-generated
plan, even an incorrect one, as a heuristic to guide greedy
best-first search, and \citet{hazra_saycanpay_2024} similarly incorporate heuristics to improve LLM-based planning. \citet{kambhampati_llms_2024} take a complementary
approach, employing critics to evaluate and refine plans produced by
LLMs. Researchers have also explored
using LLMs' code generation capabilities to produce programs that
solve specific tasks \citep{silver2024generalized}.

Prior research has also explored using LLMs to create symbolic representations of the world at runtime \cite{argenziano2024empowerembodiedmultiroleopenvocabulary}. Most closely related to our own work, \citet{musumeci2025context} propose using an LLM to generate and iteratively refine PDDL, combined with a symbolic planner to find suitable goal substitutions when the original goal is not reachable. In contrast, our approach preserves the original goal by using an LLM to reclassify available objects into the required types at runtime, enabling the original plan to remain valid without relaxing the task specification.